\documentclass[11pt,a4paper]{article}

\usepackage[margin=1in]{geometry}
\usepackage{amsmath,amssymb,amsthm}
\usepackage{graphicx}
\usepackage{booktabs}
\usepackage{tikz}
\usepackage[colorlinks=true,linkcolor=blue,citecolor=blue,urlcolor=blue]{hyperref}
\usepackage{microtype}

\newtheorem{definition}{Definition}
\newtheorem{proposition}{Proposition}

\title{Operator-Based Imaging with the DiSkO Python Package\\
\large Gridless Radio Synthesis Imaging using a Discrete Sky Operator}
\author{The DiSkO Developers\\[2mm]
\small \texttt{https://github.com/tmolteno/disko}}
\date{\today}

\newcommand{\C}{\mathbb{C}}
\newcommand{\R}{\mathbb{R}}
\newcommand{\Ns}{N_{\mathrm{s}}}
\newcommand{\Nv}{N_{\mathrm{v}}}
\newcommand{\Nb}{N_{\mathrm{b}}}
\newcommand{\Gammam}{\boldsymbol{\Gamma}}
\newcommand{\svec}{\boldsymbol{s}}
\newcommand{\vvec}{\boldsymbol{v}}

\begin{document}

\maketitle

\begin{abstract}
DiSkO (``Discrete Sky Operator'') is a Python package for synthesis imaging in
radio astronomy that avoids the traditional gridding and CLEAN/deconvolution
pipeline. Instead, the telescope is modelled directly as a finite-dimensional
linear operator $\Gammam$ that maps a pixelized sky $\svec$ to the measured
visibilities $\vvec$. The package makes the operator central: it computes the
singular value decomposition of $\Gammam$ and uses the resulting four
fundamental subspaces---the range and null spaces of the operator and of its
adjoint---to understand exactly what the telescope can and cannot measure. From
this decomposition it derives a ``natural basis'' in which the operator is
diagonal, giving closed-form images, Tikhonov-regularized images, Bayesian
sequential inference, and well-conditioned null-space-complete images. For
large problems the operator is exposed in a matrix-free form compatible with
iterative solvers (LSQR, LSMR, FISTA). This article describes the mathematics
of operator-based imaging as implemented in the package, with pointers to the
relevant classes and functions.
\end{abstract}

\section{Introduction}

In aperture-synthesis imaging, an array of antennas measures the spatial
coherence of the electric field arriving from the sky. For each pair of
antennas (a \emph{baseline}) the measured quantity---the complex
\emph{visibility}---is, under the usual assumptions of the van~Cittert--Zernike
theorem, a sample of the two-dimensional Fourier transform of the sky
brightness distribution \cite{thompson2017,perley1989}. Classical imaging
proceeds by gridding the visibilities onto a regular Fourier grid and applying
an inverse FFT followed by a deconvolution stage such as CLEAN or a
maximum-entropy method.

DiSkO takes a different route. The sky is discretized first, into a set of
$\Ns$ pixels on the celestial sphere (or on a disc-shaped field of view), and
the instrument is discretized second, into $\Nv$ visibility measurements. The
physics of the instrument is then encoded exactly, pixel by pixel, in a
measurement matrix $\Gammam\in\R^{\Nv\times\Ns}$, so that the measurement
equation is linear:

\begin{equation}
  \label{eq:measurement}
  \vvec \;=\; \Gammam\,\svec \;+\; \boldsymbol{n},
\end{equation}

where $\boldsymbol{n}$ models additive noise. There is no gridding, no
pixel-domain Fourier transform of the data, and no restoration step: the image
is obtained by (regularized) inversion of the operator $\Gammam$ itself. The
price paid is that $\Gammam$ is at least as large as the data, and in the
typical imaging regime ($\Ns > \Nv$) it is rank-deficient. The central idea of
the package is to make this rank deficiency explicit and useful by computing
the singular value decomposition (SVD) of $\Gammam$ and working in the
associated \emph{natural basis}.

This article explains how this works, following the actual implementation:
Section~\ref{sec:model} develops the discrete measurement equation;
Section~\ref{sec:svd} describes the SVD-based \texttt{TelescopeOperator} and its
four subspaces; Section~\ref{sec:natural} introduces imaging in the natural
basis; Section~\ref{sec:reg} covers regularized and sparse imaging;
Section~\ref{sec:matrixfree} describes the matrix-free operator used with
iterative solvers; Section~\ref{sec:bayes} derives the Bayesian sequential
inference used by \texttt{disko\_bayes}; Section~\ref{sec:fov} describes the
sky discretizations; and Section~\ref{sec:impl} gathers implementation notes.
Throughout, italicized code names refer to the package
\href{https://github.com/tmolteno/disko}{tmolteno/disko}.

\section{The discrete measurement equation}
\label{sec:model}

\subsection{Coordinates}

Let $(\lambda)$ be the observing wavelength. A baseline between two antennas
has coordinates $(u,v,w)$, measured in meters, in a frame where $w$ points
towards the phase centre and $(u,v)$ span the plane perpendicular to it. A
point on the sky is described by the direction cosines
$(l,m,n)$, with $l^2+m^2+n^2=1$. A unit solid angle around the direction
$(l,m,n)$ contributes a phase term
$\exp\!\big(j\,\tfrac{2\pi}{\lambda}\,(ul+vm+w(n-1))\big)$ to each visibility,
the $w(n-1)$ factor removing the geometric delay of the phase centre.

\subsection{Discretizing the sky}

A field-of-view class (Section~\ref{sec:fov}) provides $\Ns$ pixels with
direction cosines $l_j,m_j,n_j$ and quadrature weights $A_j$ (the pixel areas,
normalized to integrate to the observed solid angle). The package stores the
combination $n_j-1$ as \texttt{n\_minus\_1}. The interferometric measurement
equation then becomes the discrete sum

\begin{equation}
  \label{eq:rime}
  v_i \;=\; \sum_{j=1}^{\Ns}\;
  \underbrace{\exp\!\Big(j\,\omega\,
    \big(u_i l_j+v_i m_j+w_i(n_j-1)\big)\Big)\,A_j}_{H_{ij}}\;\;s_j,
  \qquad \omega=\frac{2\pi}{\lambda},
\end{equation}

implemented in \texttt{disko.get\_harmonics()} (vectorized with an outer
product over rows $i$ and columns $j$) and
\texttt{DiSkO.make\_gamma()}. This is a Riemann-sum discretization of the
continuous RIME \cite{hamaker1996}, and $H_{ij}$ is called a
\emph{harmonic}: it is the phase pattern that baseline $i$ records from pixel
$j$. Note the convention in the code: $s_j$ is the brightness of pixel $j$
(in units such that the pixel areas carry the solid-angle weight), and because
$\sum_j A_j^2$ is a constant row energy, every baseline row of $H$ has the same
norm (verified by \texttt{test\_harmonics} in
\texttt{disko/tests/test\_telescope\_operator.py}).

\subsection{The real-valued telescope operator}

The harmonics $H\in\C^{\Nb\times\Ns}$ are complex. Complex-valued linear
systems can be solved directly, but the package instead rewrites
\eqref{eq:rime} as a real system by stacking the real and imaginary parts of
the visibilities:

\begin{equation}
  \label{eq:real}
  \underbrace{\begin{bmatrix}\Re(\vvec)\\ \Im(\vvec)\end{bmatrix}}_{\vvec\in\R^{\Nv}}
  \;=\;
  \underbrace{\begin{bmatrix}\Re(H)\\ \Im(H)\end{bmatrix}}_{\Gammam\in\R^{\Nv\times\Ns}}
  \;\svec,
  \qquad \Nv=2\Nb,
\end{equation}

computed with a memory-efficient \texttt{np.concatenate} in
\texttt{make\_gamma()}. All subsequent operator analysis (the SVD, the natural
basis, the Bayesian updates) works with this real matrix $\Gammam$. Every
complex visibility contributes two real rows, so visibilities arriving from a
measurement set are converted with \texttt{vis\_to\_real()}.

\begin{definition}[Telescope operator]
  The \emph{telescope operator} is the matrix $\Gammam\in\R^{\Nv\times\Ns}$ of
  \eqref{eq:real}. A \emph{sky vector} is $\svec\in\R^{\Ns}$; a
  \emph{visibility vector} is $\vvec\in\R^{\Nv}$. The forward map is
  $\vvec=\Gammam\svec$.
\end{definition}

For a full-sky Healpix field of view with $n_\mathrm{side}=32$,
$\Ns=12\,n_\mathrm{side}^2=12\,288$ while a modest array produces only a few
hundred complex baselines, so $\Ns\gg\Nv$ and the operator is very
rank-deficient: the telescope simply cannot measure most of the sky.

\section{The SVD of the telescope operator}
\label{sec:svd}

The key construction of the package is the class \texttt{TelescopeOperator}
(\texttt{disko/telescope\_operator.py}). On construction it forms
$\Gammam=\texttt{grid.make\_gamma(sphere)}$ and computes the compact SVD

\begin{equation}
  \label{eq:svd}
  \Gammam \;=\; U\,\Sigma\,V^{\!\top},
\end{equation}

with $U\in\R^{\Nv\times\Nv}$, $\Sigma\in\R^{\Nv\times\Ns}$ containing the
singular values $\sigma_1\ge\sigma_2\ge\dots\ge0$ on its diagonal and $0$
elsewhere, and $V\in\R^{\Ns\times\Ns}$. (In code $V$ is the conjugate
transpose $V^{\!\top}$ returned by \texttt{scipy.linalg.svd}. Depending on a
module flag \texttt{USE\_DASK}, the factorization is computed through
\texttt{dask\_svd()} on a dask-backed array; the math is identical.)

\subsection{Rank truncation by condition number}

A telescope operator is numerically singular: measured visibilities are noisy,
and singular values below the noise floor contribute only noise amplification
if inverted. The package truncates by a condition-number budget
$\kappa=\texttt{max\_cond}=10^{4}$ (configurable): the numerical rank is
\begin{equation}
  r \;=\; \#\big\{i : \sigma_i \ge \sigma_1/\kappa\big\},
\end{equation}
and all singular values below that threshold are set to zero
(\texttt{normal\_svd}/\texttt{dask\_svd}). Equation~\eqref{eq:svd} is then a
rank-$r$ factorization. The threshold is visible in the singular-value
spectrum; Figure~\ref{fig:spectrum} shows it for the tiny synthetic telescope
used by the test-suite (4 antennas, 24 real visibility rows, a 48-pixel
Healpix sphere, rank $r=12$).

\begin{figure}[t]
  \centering
  \includegraphics[width=0.72\textwidth]{svd_spectrum.png}
  \caption{Singular-value spectrum of the real telescope operator for the
  small synthetic telescope used in the unit tests. The vertical red line
  marks the truncation: singular values to its right are treated as null
  space. (\texttt{disko/tests/test\_telescope\_operator.py}. The figure is
  reproduced by \texttt{svd\_spectrum.py} in this directory.)}
  \label{fig:spectrum}
\end{figure}

\subsection{The four fundamental subspaces}

Partition the SVD into the significant block (subscript $1$) and the
truncated block (subscript $2$):

\begin{equation}
  \Gammam \;=\;
  \begin{bmatrix}U_1 & U_2\end{bmatrix}
  \begin{bmatrix}\Sigma_1 & 0\\ 0 & 0\end{bmatrix}
  \begin{bmatrix}V_1^{\!\top}\\ V_2^{\!\top}\end{bmatrix},
\end{equation}

with $U_1\in\R^{\Nv\times r}$, $U_2\in\R^{\Nv\times(\Nv-r)}$,
$V_1\in\R^{\Ns\times r}$, $V_2\in\R^{\Ns\times(\Ns-r)}$, and diagonal
$\Sigma_1=\operatorname{diag}(\sigma_1,\dots,\sigma_r)$. Because $U$ and $V$
are orthogonal, the columns of these four blocks are orthonormal bases of four
mutually orthogonal subspaces, which completely describe the instrument:
\cite{nullspace2016}

\begin{center}
\begin{tabular}{lll}
\toprule
Basis & Subspace & Meaning \\
\midrule
$U_1$ & range of $\Gammam$, $\R^{\Nv}$ & visibilities the telescope \\
      & $=\mathcal{R}(\Gammam)$ & can actually produce \\
$U_2$ & null space of $\Gammam^{\!\top}$ & ``invisible'' data directions \\
      & $=\mathcal{N}(\Gammam^{\!\top})$ & (never excited by any sky) \\
$V_1$ & range of $\Gammam^{\!\top}$, $\R^{\Ns}$ & measurable sky \\
      & $=\mathcal{R}(\Gammam^{\!\top})$ & (``range harmonics'') \\
$V_2$ & null space of $\Gammam$ & invisible sky \\
      & $=\mathcal{N}(\Gammam)$ & (``null harmonics'') \\
\bottomrule
\end{tabular}
\end{center}

In particular the sky space decomposes as the orthogonal direct sum
\begin{equation}
  \R^{\Ns} \;=\; \mathcal{R}(\Gammam^{\!\top}) \oplus \mathcal{N}(\Gammam),
\end{equation}
i.e.\ every sky is the sum of a \emph{measurable} part (spanned by $V_1$) and
an \emph{invisible} part (spanned by $V_2$) that produces exactly zero
visibilities. The four blocks are exposed by the API:
\texttt{U\_1}, \texttt{U\_2}, \texttt{V\_1}, \texttt{V\_2},
\texttt{rank} ($r$), \texttt{n\_r()}, \texttt{n\_n()}; the basis vectors are
retrievable via \texttt{range\_harmonic(h)} and \texttt{null\_harmonic(h)}.

\subsection{Projections}

Orthogonal projectors follow immediately from the orthonormal bases:

\begin{equation}
  P_r = V_1V_1^{\!\top}
  \quad(\text{onto measurable sky, \texttt{P\_r()}}),\qquad
  P_n = V_2V_2^{\!\top}
  \quad(\text{onto invisible sky, \texttt{sky\_to\_null}}).
\end{equation}

The package deliberately never forms the dense $P_n$ (an
$\Ns\times\Ns$ matrix); \texttt{sky\_to\_null} applies the two factors in
sequence, which costs $\mathcal{O}(\Ns(\Ns-r))$ instead of
$\mathcal{O}(\Ns^2)$ storage.

\subsection{The natural basis}

Since $V$ is orthogonal, the change of coordinates
\begin{equation}
  \label{eq:natural}
  \svec = V\boldsymbol{x}
  \qquad\Longleftrightarrow\qquad
  \boldsymbol{x}=V^{\!\top}\svec
\end{equation}
rotates the sky into the \emph{natural basis}
(\texttt{natural\_to\_sky}/\texttt{sky\_to\_natural}). In this basis the
forward model becomes diagonal:

\begin{equation}
  \vvec = \Gammam\svec = U\Sigma V^{\!\top} V\boldsymbol{x}
       = \underbrace{U\Sigma}_{\text{\texttt{A}}}\boldsymbol{x},
  \qquad
  \text{\texttt{A}} =
  \begin{bmatrix} A_r & 0\\ 0 & 0\end{bmatrix},
  \qquad A_r = U_1\Sigma_1 \in \R^{\Nv\times r}.
\end{equation}

The first $r$ components $x_r=V_1^{\!\top}\svec$ are the measurable
coefficients; the last $\Ns-r$ components $x_n=V_2^{\!\top}\svec$ are
invisible. Because $U_1^{\!\top}U_1=I_r$ and $\Sigma_1$ is diagonal, the
truncated operator $A_r$ is diagonal in the natural basis up to an orthogonal
rotation of the data: for a sky in the range space, the data is
$\vvec = U_1\Sigma_1 x_r$ and therefore

\begin{equation}
  \label{eq:diag}
  U_1^{\!\top}\vvec \;=\; \Sigma_1 x_r,
\end{equation}

a system that is inverted by division on the diagonal. All quantitative imaging
and inference in the package is built on this observation.

\section{Imaging in the natural basis}
\label{sec:natural}

\subsection{Direct (least-squares) imaging}

The most direct reconstruction ignores the decomposition and solves
\eqref{eq:measurement} in the least-squares sense:
\texttt{image\_visibilities} computes
$\hat{\svec}=\Gammam^{+}\vvec$ via \texttt{np.linalg.lstsq}. This is the
minimum-norm image, but it can be numerically unstable and mixes the measurable
content with poorly-conditioned modes.

\subsection{Natural-basis imaging}

Better use is made of the structure. \texttt{image\_natural} implements the
closed form derived from \eqref{eq:diag}:

\begin{enumerate}
  \item project the data onto the range space: $\boldsymbol{y}=U_1^{\!\top}\vvec$
        (\texttt{U\_1.T @ vis});
  \item undo the diagonal: $x_r=\Sigma_1^{-1}\boldsymbol{y}$, an
        element-wise division by the singular values;
  \item rotate back: $\hat{\svec}=V_1 x_r$
        (\texttt{V\_1 @ x\_r}).
\end{enumerate}

This is imaging by explicit inversion of the well-conditioned part of the
instrument: the reconstruction is the sky component that the telescope can
measure, recovered with the correct amplitude and phase of every mode.

\subsection{Null-space completion}

The reconstruction of Section~\ref{sec:natural} leaves the null-space
component of the image exactly zero. If an external prior sky
$\svec_{\mathrm{prior}}$ is available (e.g.\ a map from another survey),
\texttt{image\_natural(null\_prior=...)} / \texttt{complete\_image} grafts its
invisible component onto the reconstruction:

\begin{equation}
  \label{eq:complete}
  \hat{\svec}_{\mathrm{complete}}
  \;=\; V_1\Sigma_1^{-1}U_1^{\!\top}\vvec
     \;+\; V_2V_2^{\!\top}\svec_{\mathrm{prior}}.
\end{equation}

By construction $\Gammam V_2=0$, so the added term contributes \emph{zero}
visibilities: the completed image fits the measured data \emph{exactly}, while
carrying the prior's structure in the directions the telescope cannot see
(verified in \texttt{test\_complete\_image}). The \texttt{disko\_bayes}
command-line tool exposes this as \texttt{--null-prior}.

\section{Regularized imaging}
\label{sec:reg}

When $\Ns>\Nv$ the least-squares problem is under-determined and the null
space must be handled. The package offers two classical regularizations in
addition to the null-space completion above.

\subsection{Tikhonov (L2) regularization}

\texttt{image\_tikhonov} computes the minimum of
$\|\vvec-\Gammam\svec\|^2+\alpha^2\|\svec\|^2$. Using the SVD this has the
famous closed form; the package exploits the truncation to apply the filter
only on the rank block:

\begin{equation}
  \label{eq:tikhonov}
  \hat{\svec}_\alpha
  \;=\; V_1\,
    \Big[ \tfrac{\sigma_i}{\sigma_i^2+\alpha^2}\Big]_{i=1}^{r}\,
    \big(U_1^{\!\top}\vvec\big),
\end{equation}

i.e.\ each singular value $\sigma_i$ is attenuated by the factor
$\sigma_i/(\sigma_i^2+\alpha^2)$ while the null space stays at zero. Since
singular values below the rank cutoff are already zero, there is nothing to
regularize outside the rank block; working with the $V_1,U_1$ factors is
$\mathcal{O}(\Ns r)$ rather than a dense
$\mathcal{O}(\Ns^2)$ triple product
(\texttt{CHANGES.md}, Unreleased).

\subsection{Sparsity-enforcing (L1) imaging}

For sources that are sparse on the sky, \texttt{image\_lasso} solves the
elastic-net problem
\begin{equation}
  \min_{\svec}\; \tfrac12\|\vvec-\Gammam\svec\|^2
  \;+\; \alpha\Big(\rho\|\svec\|_1 + \tfrac{1-\rho}{2}\|\svec\|_2^2\Big),
\end{equation}
via \texttt{sklearn.linear\_model.ElasticNet} (optionally
\texttt{ElasticNetCV} with cross-validated $\alpha$). The sparsity of the L1
term suppresses the null-space leakage that plain least squares would
introduce, at the cost of a non-linear solve. The regularization strength is
scaled as $\alpha/\sqrt{\Ns}$ to keep the norm of the L1 penalty comparable
across pixel counts.

\section{Matrix-free operators and iterative solvers}
\label{sec:matrixfree}

Explicitly forming and factorizing $\Gammam$ (as
Section~\ref{sec:svd} does) costs $\mathcal{O}(\Ns\Nv)$ memory and is
impossible for large telescopes or dense skies. For the iterative and sparse
paths the package instead exposes a \emph{matrix-free} operator,
\texttt{DiSkOOperator} (\texttt{disko/disko.py}), a subclass of
\texttt{pylops.LinearOperator} \cite{pylops2018}.

\subsection{The operator \texttt{DiSkOOperator}}

The operator has shape $M\times N$ with $M=2\Nb\cdot N_{\mathrm{freq}}$ (real
and imaginary rows, per frequency) and $N=\Ns$. Instead of storing $\Gammam$
it implements

\begin{itemize}
  \item \texttt{\_matvec(x)}: computes
    $\vvec = \begin{bmatrix}\Re(H)\\ \Im(H)\end{bmatrix}\svec$ by summing, over
    frequencies, the products of the real and imaginary parts of $H$ with the
    sky. The harmonics are evaluated in \emph{blocks} of visibility rows
    (\texttt{\_block\_size()} targets $\sim$200\,MB per block) and cached as
    float32 arrays, valid because the UVW coordinates and frequencies are
    constant across iterations;
  \item \texttt{\_rmatvec(v)}: the exact adjoint
    $\svec = \Re(H)^{\!\top}\vvec_{\mathrm{re}}+
    \Im(H)^{\!\top}\vvec_{\mathrm{im}}$, the so-called dirty image (the
    gridless primary beam).
\end{itemize}

Both are linear in the data and agree with the dense $\Gammam$
(\texttt{test\_pylops\_operator.py} checks element-by-element agreement and
passes the pylops dot-test). The operator flag \texttt{explicit=False}
indicates that it can only be applied, never inverted directly.

\subsection{Iterative solvers}

\texttt{DiSkO.solve\_matrix\_free} plugs this operator into three solvers:

\begin{itemize}
  \item \textbf{LSQR} \cite{paige1982} via
    \texttt{scipy.sparse.linalg.lsqr}, with damping parameter
    \texttt{alpha} ($\ge0$; negative selects the data RMS);
  \item damped least squares on the normal equations (the
    \texttt{--lsmr} path) via
    \texttt{pylops.optimization.leastsquares.normal\_equations\_inversion};
  \item \textbf{FISTA} \cite{beck2009} via
    \texttt{pylops.optimization.sparsity.fista} for sparse (L1) imaging, using
    the soft-threshold form and initialized with the true adjoint image
    $|\texttt{A.T} @ \texttt{d}|$ so that the starting guess is physically
    non-negative and correctly oriented (a fix recorded in v1.4.1).
\end{itemize}

Because each iteration needs only matvec/rmatvec, memory use is dominated by
the cached harmonic blocks, not by $\Gammam$ itself. This same operator backs
the plain dirty image used by \texttt{image\_visibilities} in \texttt{DiSkO}
(adjoint applied with blocked matrix-free code).

\section{Bayesian sequential inference}
\label{sec:bayes}

The most sophisticated use of the operator structure is full Bayesian
inference, implemented in \texttt{bayes\_cli.py} and
\texttt{TelescopeOperator.sequential\_inference}. Let the prior on the sky be a
multivariate Gaussian

\begin{equation}
  p(\svec)\;=\;\mathcal{N}(\svec\mid \boldsymbol{\mu}_0,\Sigma_0),
\end{equation}

with the heuristic default
$\boldsymbol{\mu}_0=p_{50}(|\vvec|)\mathbf{1}$ (median visibility amplitude)
and $\Sigma_0=p_{95}(|\vvec|)^2 I$ (\texttt{get\_prior},
\texttt{create\_prior}). The likelihood for the measured (real) visibilities is
$p(\vvec\mid\svec)=\mathcal{N}(\vvec\mid\Gammam\svec,\Sigma_v)$ with a
visibility covariance built from the per-baseline RMS.

\subsection{The update}

The natural basis makes the computation tractable. With
$\boldsymbol{x}=V^{\!\top}\svec$ and $\vvec=A\boldsymbol{x}$, the prior is
rotated into the natural basis and split into its range block
($x_r$, size $r$) and null block ($x_n$, size $\Ns-r$). The likelihood
couples only to $x_r$ through $A_r=U_1\Sigma_1$
(\texttt{prior\_r.bayes\_update(precision, vis, A\_r)}), which is a standard
Gaussian update \cite{bishop2006}:

\begin{align}
  \Sigma_r^{-1} &= \Sigma_{r,0}^{-1} + A_r^{\!\top}\Sigma_v^{-1} A_r,\\
  \boldsymbol{\mu}_r &= \Sigma_r \big(\Sigma_{r,0}^{-1}\boldsymbol{\mu}_{r,0}
  + A_r^{\!\top}\Sigma_v^{-1}\vvec\big),
\end{align}

while the null block receives \emph{no} information from the data:
$\Sigma_n=\Sigma_{n,0}$, $\boldsymbol{\mu}_n=\boldsymbol{\mu}_{n,0}$. The
posterior is the Cartesian product of the two blocks, rotated back to the sky
space with $V$ (\texttt{MultivariateGaussian.outer} then
\texttt{linear\_transform(V)}).

Two exact optimizations are used:

\begin{itemize}
  \item \textbf{Diagonal prior fast path}: when
    $\Sigma_0=\sigma_0^2 I$ (detected by
    \texttt{is\_scaled\_identity}), $V^{\!\top}\Sigma_0 V=\sigma_0^2 I$
    stays diagonal, so only the mean needs rotating, and the posterior is
    assembled directly as
    $\Sigma = V_1\Sigma_r V_1^{\!\top}+ \sigma_0^2(I-V_1V_1^{\!\top})$,
    avoiding two $\mathcal{O}(\Ns^3)$ matrix products
    (\texttt{do\_inference} in \texttt{bayes\_cli.py});
  \item the visibility precision $\Sigma_v^{-1}$ is formed once
    (\texttt{MultivariateGaussian.sp\_inv}).
\end{itemize}

\subsection{Sequential (multi-snapshot) inference}

Because the posterior is again a multivariate Gaussian, it can be fed back as
the prior of the next snapshot. The \texttt{--hdf} mode of
\texttt{disko\_bayes} does exactly this, looping over visibilities and chaining
\texttt{posterior $\to$ prior} through the single-update machinery above
(asserted in \texttt{test\_do\_inference\_sequential}).
A non-diagonal starting prior (e.g.\ a posterior loaded with
\texttt{--prior}) takes the same dense per-turn chaining. Because each
snapshot has its own natural basis, the null block is carried forward in that
snapshot's basis; this is exact within a single update, and across turns only
when the successive range spaces coincide.

For the dedicated \texttt{--sequential N} option (N turns, each drawing
\texttt{nvis} fresh visibilities from the measurement set) the chain instead
runs in a low-rank \emph{information form} (class
\texttt{SequentialInfoState}) that never leaves the sky frame.
With a scaled-identity prior $\Sigma_0=\sigma_0^2 I$, the posterior after $k$
turns has the precision
\begin{equation}
  \label{eq:infoprec}
  \Sigma_k^{-1}
  \;=\; \sigma_0^{-2}I + \sum_{j=1}^{k} V_{1,j} B_j V_{1,j}^{\!\top},
  \qquad
  w_k
  \;=\; \sigma_0^{-2}\boldsymbol{\mu}_0 + \sum_{j=1}^{k} V_{1,j} c_j,
\end{equation}
with $B_j=A_{r,j}^{\!\top}\Sigma_{v,j}^{-1}A_{r,j}$ and
$c_j=A_{r,j}^{\!\top}\Sigma_{v,j}^{-1}\vvec_j$ the range-block precision and
data terms of turn $j$. The chain stores only the stacked range bases
$L_k=[V_{1,1}\,\dots\,V_{1,k}]\in\R^{\Ns\times R}$ ($R=\sum_j r_j$ the
cumulative rank) and the compressed inverse
$M_k=\big(C_k^{-1}+\sigma_0^2 L_k^{\!\top}L_k\big)^{-1}$ with
$C_k=\operatorname{blockdiag}(B_1,\dots,B_k)$, maintained by a bordered Schur
update per turn; no $\Ns\times\Ns$ matrix is ever formed and no covariance is
rotated between bases. Sky-space quantities are recovered with Woodbury
identities only when output is needed:
\begin{equation}
  \label{eq:woodbury}
  \mu_k=\boldsymbol{\mu}_0+\sigma_0^2 L_k M_k
    \big(C_k^{-1}c_k-L_k^{\!\top}\boldsymbol{\mu}_0\big),
  \qquad
  \Sigma_k=\sigma_0^2 I-\sigma_0^4 L_k M_k L_k^{\!\top},
\end{equation}
so the per-step mean image costs $\mathcal{O}(R^2+R\Ns)$, the per-pixel
variance $\mathcal{O}(R^2\Ns)$ (the diagonal of the low-rank product), and the
point-covariance row comes from $M_kL_k^{\!\top}$ without materializing
$\Sigma_k$. This is the exact conjugate-Gaussian posterior of
\eqref{eq:infoprec} for the sequence of reduced operators---no information is
dropped between turns---and is asserted against a direct precision-form
construction in \texttt{test\_sequential.py}, together with the identity
$\Sigma_k\Sigma_k^{-1}=I$ and the property that each turn's data adds
precision only along that turn's own range space. The package also draws
posterior samples (with results written as \texttt{*\_mu}, \texttt{*\_var},
\texttt{*\_pcf} and \texttt{*\_sN} images) and reports the point-covariance
function at the brightest pixel.

\subsection{Variance shrinkage, localization, and conditioning}
\label{sec:variance}

Carrying the posterior in information form makes two structural properties of
the sequential posterior explicit. Both hold exactly, and both are exercised
step by step by the unit tests.

\begin{proposition}[Monotone variance shrinkage]
\label{prop:shrink}
Each turn adds a positive semidefinite term to the precision
\eqref{eq:infoprec}, so the posterior covariances decrease in the Loewner
order,
\begin{equation}
  \Sigma_{k+1}^{-1} \;=\; \Sigma_k^{-1} + V_{1,k+1}B_{k+1}V_{1,k+1}^{\!\top}
  \;\succeq\; \Sigma_k^{-1}
  \qquad\Longrightarrow\qquad
  \Sigma_{k+1} \;\preceq\; \Sigma_k .
\end{equation}
Every pixel's variance therefore weakly decreases from turn to turn
(diagonals obey the same ordering, since
$e_i^{\!\top}\Sigma_{k+1}e_i\le e_i^{\!\top}\Sigma_ke_i$), the decrease is
strict at every pixel with a nonzero component along the new turn's range
basis $V_{1,k+1}$ (for which $e_i^{\!\top}V_{1,k+1}B_{k+1}V_{1,k+1}^{\!\top}
e_i>0$), and the trace falls monotonically below its prior value
$\Ns\sigma_0^2$.
\end{proposition}

\begin{proposition}[Null-space localization]
\label{prop:localize}
Let $v\in\mathcal{N}(L_k^{\!\top})$ be a sky direction orthogonal to every
absorbed range basis. The Woodbury form \eqref{eq:woodbury} then gives
$v^{\!\top}\Sigma_k v=\sigma_0^2\|v\|^2$ and
$v^{\!\top}(\boldsymbol{\mu}_k-\boldsymbol{\mu}_0)=0$: such directions keep
the prior variance and prior mean \emph{exactly}, at every step of the chain.
In particular, a pixel whose basis vector $e_i$ has zero entries in all
stacked range bases $V_{1,1},\dots,V_{1,k}$---a pixel no turn has ever
projected onto---is never moved by the data.
\end{proposition}

Two consequences deserve note. First, if successive range bases have
\emph{disjoint supports} (so that
$\operatorname{span}(V_{1,i})\perp\operatorname{span}(V_{1,j})$ for $i\neq j$),
the precision sum decomposes over orthogonal subspaces and the posterior is
block diagonal: turn $k+1$ then changes neither the variance nor the mean of
pixels informed only by earlier turns---data cannot perturb a posterior
coordinate it does not project onto. Second, for real telescope geometries the
range bases are generic and no pixel direction is exactly orthogonal to the
span of all of them: \emph{every} pixel acquires at least a little
information, so all per-pixel variances decrease at every step, and exact
null pixels arise only by construction. Both propositions are verified against
hand-built range bases with exact zero rows
(\texttt{TestVarianceLocalization} in \texttt{test\_sequential.py}), the
monotone decrease is asserted at every step of the chain on distinct turns,
and the chain is also exercised end-to-end on synthetic data generated with
the package's own forward operator: the tartball tool \cite{tartball}
predicts a known spot model through the operator for the 24-antenna TART
layout onto an $n_\mathrm{side}=16$ healpix sphere, adds Gaussian noise, and
writes a measurement set (\texttt{test\_data/synth.ms}); the test draws
disjoint pools
of baselines from it exactly as \texttt{run\_sequential} does, and verifies
Proposition~\ref{prop:shrink} at every step together with a data-fit
improvement of the posterior mean over the prior mean
(\texttt{TestSequentialSyntheticMs}).

\emph{Numerical conditioning of the covariance accessors.} The sky-space
covariance quantities evaluate the Woodbury difference
$\sigma_0^2-\sigma_0^4\operatorname{diag}(L_kM_kL_k^{\!\top})$, a
cancellation of comparable terms whose relative error grows with the
likelihood-to-prior precision ratio. On the 48-pixel test problem the
per-pixel variance recovered by \texttt{variance()} agrees with a dense
precision-form oracle to machine precision when
$\Sigma_v^{-1}\sim\sigma_0^{-2}$, to $4\times10^{-9}$ at
$\sigma_v=10^{-4}$ and $4\times10^{-5}$ at $\sigma_v=10^{-6}$; beyond that
($\sigma_v\lesssim10^{-8}$) the recovered sky-space covariance is dominated by
round-off, although the internal information state $(\boldsymbol{w}_k, M_k,
L_k)$ remains accurate---only the back-transfer to sky space degrades (and a
dense oracle degrades equally, since
$\operatorname{cond}(\Sigma_k^{-1})$ grows with the same ratio). The
well-conditioned operating point is therefore a realistic per-baseline noise
level with likelihood and prior precisions comparable---the regime of the
\texttt{--sigma-v} values used by \texttt{disko\_bayes} (e.g.\
$\sigma_v=0.15$)---and the exactness tests use such well-conditioned
likelihoods for all covariance comparisons.

\section{Discretizations of the sky}
\label{sec:fov}

All the machinery above is agnostic to the sky discretization; a field-of-view
class need only provide \texttt{npix}, \texttt{l}, \texttt{m}, \texttt{n},
\texttt{n\_minus\_1} and \texttt{pixel\_areas}. The package supplies:

\begin{itemize}
  \item \texttt{HealpixFoV}: the full celestial sphere with a HEALPix
        tessellation \cite{gorski2005}, equal-area pixels, used for wide-field
        or all-sky imaging;
  \item \texttt{HealpixSubFoV}: a disc-shaped subset of a HEALPix grid
        centred on a phase centre (built with \texttt{healpy.query\_disc}),
        with pixel areas renormalized to the disc;
  \item \texttt{SquareFoV}: a flat-sky square with a fixed resolution in
        arcminutes, useful for small fields where curvature is negligible;
  \item \texttt{AdaptiveMeshFoV}: an unstructured triangular mesh generated
        with \texttt{gmsh} whose resolution can locally refine
        (\texttt{refine()}) in regions of high gradient during repeated
        Tikhonov imaging (the \texttt{--mesh} / \texttt{--adaptive} CLI
        options).
\end{itemize}

The discretization is what makes the approach \emph{gridless} in one sense:
pixels may be anywhere on the sphere, at arbitrary sizes, and the operator
accounts for each pixel's solid angle exactly. The pixel areas $A_j$ appear in
the harmonics \eqref{eq:rime}, so imaging with any of these skies automatically
handles varying pixel size.

\section{Implementation notes}
\label{sec:impl}

A few implementation choices are worth highlighting because they shape the
mathematics.

\begin{itemize}
  \item \textbf{Dask}: for large operators, $\Gammam$ is wrapped in a dask
        array (\texttt{USE\_DASK}) and the SVD is run on it
        (\texttt{dask\_svd}); all basis methods (\texttt{sky\_to\_null},
        \texttt{natural\_to\_sky}, \texttt{image\_natural}, \dots) force the
        lazy graphs eagerly and return plain NumPy arrays
        (\texttt{test\_explicit\_computation}).
  \item \textbf{SVD cache}: the factorization is expensive, so it is cached to
        disk (\texttt{svd\_Ns\_Nv\_maxcond.npz}) keyed by the problem size and
        the truncation factor, and reloaded when available
        (\texttt{use\_cache=True}).
  \item \textbf{Conditioning}: the rank cutoff \texttt{max\_cond} is exposed
        through \texttt{TelescopeOperator(..., max\_cond=...)} and
        \texttt{disko\_bayes --max-cond}; a smaller budget moves more modes
        into the null space (checked in \texttt{test\_max\_cond}).
  \item \textbf{Data conventions}: measurement-set readers restore the
        conjugate-symmetric visibility pairs ($\pm(u,v,w)$), so the operator
        rows come in conjugate pairs; visibilities are natural-weight divided
        by the per-baseline RMS before imaging, and \texttt{disko} conjugates
        them where the UVW convention in the data file requires it.
  \item \textbf{Complexity}: the dense SVD is the main cost,
        $\mathcal{O}(\min(\Ns\Nv^2,\Nv\Ns^2))$; it is paid once. Imaging in
        the natural basis is then $\mathcal{O}(\Ns r)$; Tikhonov imaging adds
        only the diagonal filter; Bayesian updates cost
        $\mathcal{O}(r^3)$ for the range-block inversion plus
        $\mathcal{O}(\Ns r)$ rotations.
\end{itemize}

\section{Conclusion}

Operator-based imaging in DiSkO replaces the classic grid--FFT--CLEAN pipeline
with a single linear-algebra object: the telescope operator $\Gammam$ and its
SVD. The SVD reveals exactly what the telescope measures (the range space
$V_1$), what it cannot measure (the null space $V_2$), and provides a natural
basis in which the operator is diagonal. On this basis the package builds
closed-form images, Tikhonov and sparse regularizations, exactly
data-consistent null-space-completed images, and full Bayesian inference with
analytic sequential updates whose uncertainty shrinks monotonically and
exactly on the subspace the data informs. For problems too large for an
explicit SVD, the
same operator is exposed matrix-free to LSQR/LSMR/FISTA. The result is an
imaging pipeline whose properties---resolution, information content,
regularization, uncertainty---are all read directly off the operator rather
than from ad-hoc deconvolution heuristics.

\begin{thebibliography}{9}

\bibitem{thompson2017}
A.~R. Thompson, J.~M. Moran, and G.~W. Swenson.
\emph{Interferometry and Synthesis in Radio Astronomy}, 3rd ed.
Springer, 2017.

\bibitem{perley1989}
R.~A. Perley, F.~R. Schwab, and A.~H. Bridle (eds.).
\emph{Synthesis Imaging in Radio Astronomy}.
ASP Conference Series, vol.~6, 1989.

\bibitem{hamaker1996}
J.~P. Hamaker, J.~D. Bregman, and R.~J. Sault.
Understanding radio polarimetry. I. Mathematical foundations.
\emph{Astron.\ Astrophys.\ Suppl.\ Ser.}, 117:137--147, 1996.

\bibitem{nullspace2016}
Null space from SVD, Mathematics Stack Exchange,
\url{https://math.stackexchange.com/questions/1771013/null-space-from-svd}.

\bibitem{paige1982}
C.~C. Paige and M.~A. Saunders.
LSQR: An algorithm for sparse linear equations and sparse least squares.
\emph{ACM Trans.\ Math.\ Softw.}, 8(1):43--71, 1982.

\bibitem{beck2009}
A.~Beck and M.~Teboulle.
A fast iterative shrinkage-thresholding algorithm for linear inverse problems.
\emph{SIAM J.\ Imaging Sci.}, 2(1):183--202, 2009.

\bibitem{bishop2006}
C.~M. Bishop.
\emph{Pattern Recognition and Machine Learning}.
Springer, 2006 (Section~2.3, Bayesian inference for the Gaussian).

\bibitem{gorski2005}
K.~M. G\'orski, E.~Hivon, A.~J. Banday, B.~D. Wandelt, F.~K. Hansen,
M.~Reinecke, and M.~Bartelmann.
HEALPix: A framework for high-resolution discretization and fast analysis of
data distributed on the sphere.
\emph{Astrophys.\ J.}, 622:759--771, 2005.

\bibitem{pylops2018}
PyLops: A linear-operator library for Python,
\url{https://pylops.readthedocs.io}.

\bibitem{disko}
T.~Molteno et al., DiSkO: Discrete Sky Operator synthesis imaging,
\url{https://github.com/tmolteno/disko}.

\bibitem{tartball}
T.~Molteno et al., tartball: simulation of TART radio telescope data,
\url{https://github.com/tart-telescope/tartball}.

\end{thebibliography}

\end{document}
